\documentclass[11pt]{article}
\usepackage[a4paper,margin=24mm]{geometry}
\usepackage{amsmath,amssymb,amsthm,hyperref}
\newtheorem{theorem}{Theorem}\newtheorem{lemma}{Lemma}
\title{Uniform complexity of the pure-cell insertion construction}
\author{Complete proof; three independent mathematical reviews completed}
\date{30 September 2026}
\begin{document}\maketitle
This note audits the dependence on the number of controlled moments in
the fixed-number construction. It bounds the data that determine the
\emph{small} dice. Denominators of higher polynomial coefficients,
Riesz solves, and the final large dice are deliberately not bounded.

Take the complexity parameter $P\ge2$ to be an integer, replacing it
by its ceiling when necessary. Call a pure-cell model of complexity at
most $P$ if the following
quantities are at most $P$: the number of labels and cells, the common
endpoint-grid denominator, the common denominator of all prescribed
raw cell moments, polynomial degrees, and the maximum and reciprocal
positive minimum of each active density piece. Ordered target degree
and the number of prescribed low moments are included among these
quantities. Each label has total mass one. An inactive density is zero
and is excluded from the positive-minimum condition.

\begin{lemma}[Uniform local parameter dependence]\label{local}
There is an absolute constant $C$ such that the synchronized local
completion theorem can be carried out with all its controlled output
data bounded by
\[
 P^{(q+2)^C}.
\]
Here the input family has at most $P$ instances, degree and target
order at most $P$, density bounds $P^{-1}\le w_\alpha\le P$,
$q\le P$, prescribed node counts $K_\alpha\le P$, and the quantities
$K_\alpha\int t^j w_\alpha$, $j\le2q+1$, have one common denominator
at most $P$. The counts must be sufficiently large for the separated
quadrature; a sufficient lower bound has the form
\[
 K_\alpha\ge e^{C(q+1)}
    (a^{-1}+b+1)^C(m+\deg w_\alpha+q+2)^4,
\]
with a change of the absolute constant, where $a,b$ are the actual
common density bounds. Increasing the fixed polynomial degree in this
last display would make no difference to any conclusion below.
\end{lemma}
\begin{proof}
We give explicit uniform types of all losses in
\path{synchronized_local_moment_completion.tex}. This avoids treating
its fixed-$q$ constants as unspecified functions.
For this audit, a bound is of the permitted type if its logarithm is
at most $(q+2)^C\log P$ for some absolute $C$. A bounded number of
products, fixed powers, or powers of order polynomial in $q$ preserve
this class after increasing $C$.

The separated weighted quantile/Newton lemma has only fixed powers of
$a^{-1},b,m+\deg w+q+2$. For example
$\|Q'\|_\infty\le a^{-1}$ and
$\|Q''\|_\infty\le C(\deg w+1)^2ba^{-3}$ for the quantile map.
The derivative Gram is bounded below by $a$, and the polynomial
Markov bounds for its perturbation have fixed powers. Forcing $q+2$
control nodes costs a further polynomial in $q$: delete $q+2$ Gram
columns and move that many nodes by $O(1/K)$. The required quadrature
degree is at most $C3^q(m+\deg w+q+2)$. These observations give the
stated sufficient lower bound after enlarging $C$.

The universal feature conditioning is now quantitative in
\path{uniform_prefix_conditioning.tex}. For perturbation size $\eta$,
its inverse norm is at most
\[
 P^{C(q+1)^3}\eta^{-2q}.
\]
The inverse Gram and the fixed Newton right inverse therefore cost
at most another fixed power of this quantity, times fixed powers of
$P$ and $q$. Norms of all feature derivatives are bounded by fixed
powers of their degrees, with factors at most polynomial in $q$.
Thus, once $\eta^{-1}$ is of the permitted type, sufficiently fine
rational rank rounding preserves the fixed-inverse contraction with
permitted precision. There is no product over the number of Newton
iterations: one fixed right inverse and its contraction-radius bound
control the complete solve.

For clarity, the remaining uniform losses are as follows.
\begin{enumerate}
\item Rescaled moment systems on a gap or box, with orders at most
$q+1$, have inverse norms bounded by
$e^{C(q+1)^2\log(q+2)}\Delta^{-C(q+1)^2}$, where $\Delta$ is the
cell length. This follows directly from the triangular Legendre
basis, binomial coefficients, and the nonzero diagonal moments.
\item The degree-$q$ smoothed quotient estimate has inverse powers
of the minimum gap of order $q+3/2$ and polynomial degree power
$2q+4$. Its constants can be bounded by $e^{C(q+1)^2\log(q+2)}$.
Indeed its proof uses $q+1$ derivatives, repeated integration, and
binomial jet cancellation. The inverse Markov bound is
$(C M^2/\Delta)^{q+1}$. The auxiliary backward recurrence over the
polynomial degree is a geometric series with ratio at most $1/2$,
not a product of degree-many constants.
\item The shared control-feature matrix has size $q+2$. Its baseline
is the rational Vandermonde matrix at $i/(q+3)$, with rows
$1,t,t^2/2,\ldots,t^{q+1}/(q+1)$. Its inverse norm is at most
$e^{C(q+1)^2\log(q+2)}$ by its determinant and cofactors. If its
rounded entries have denominator $L$, a common denominator for its
rational inverse is at most $(q+2)!L^{q+2}$, up to a fixed power of
$L$; all entries are bounded. Thus the same inverse used in every
instance introduces a permitted common denominator.
\item The low Legendre modes for box moments cost at most
$e^{C(q+1)^2\log(q+2)}\epsilon^{-C(q+1)^2}$.
The next mode, of degree $q+1$, has leading response
$\epsilon^{q+1}c_q$ times the already controlled feature-evaluation
matrix, where $|c_q|^{-1}\le e^{C(q+1)\log(q+2)}$.
Taylor remainder bounds have degree losses of order $D^{O(q)}$ and
powers of $\epsilon$ of order polynomial in $q$. Thus restoring higher moments preserves all
low conditional moments at permitted precision.
\end{enumerate}
The constants in the first two items can alternatively be obtained
by expanding the explicit proofs in
\path{higher_prefix_positive_completion.tex} and
\path{../independent_directions/quantitative_smoothed_quotient.tex};
no matrix determinant of size comparable to $m$ is used to bound a
controlled denominator.

The gap lower bound before smoothing is an inverse fixed power of
$P$. Choose $\eta=P^{-(q+2)^{C_1}}$ to meet the positive-completion
bound with slack, then a common box width
$\epsilon=P^{-(q+2)^{C_2}}$, with $C_2$ large enough for the
conditioned smoothing and the box-measure error in the Riesz solve.
Finally choose the feature-rounding, center-rounding, and low-moment
rounding grids in that order with exponents $(q+2)^{C_3}$,
$(q+2)^{C_4}$, $(q+2)^{C_5}$, increasing these absolute exponents as
needed. They can be chosen as integer powers of an integer at most
$2P$, with the prescribed common arithmetic factors multiplied in.
Those factors have permitted size: the input denominator, $(q+1)!$,
and the one shared control minor just discussed. This is a fixed
number of choices, independent of $m$ and $q$.

Every smallness requirement in the reviewed construction is a product
of the losses just listed. Hence these finitely many exponents can
all be dominated by one absolute exponent $C$.
High-dimensional rational minimum-norm or Riesz solves may create
large coefficient denominators, but their numerical norms are bounded
by the displayed coercivity/quotient estimates. Their prescribed low
moments and their already fixed endpoints do not change. Positivity,
cell lengths, masses, and degrees therefore retain permitted bounds.
This proves the lemma.
\end{proof}

\begin{lemma}[Uniform grid reset]\label{reset}
Under the complexity convention above, the uniform-grid moment reset
with reserve $q$ and ordered target degree at most $P$ produces
complexity at most $P^{(q+2)^C}$, for an absolute $C$.
\end{lemma}
\begin{proof}
In \path{uniform_grid_moment_refinement.tex}, choose a grid refining
the old endpoint grid by a factor at most $q+1$, increasing it by
another fixed power of $P$ if needed. The translated adjacent-cell
moment determinant is
\[
 Q^{-(q+1)(q+2)/2}
       \det\left(\int_j^{j+1}y^rdy\right)_{r,j=0}^q.
\]
The fixed determinant is the Vandermonde determinant at
$0,\ldots,q$, hence a positive integer $\prod_{j=1}^q j!$.
Entries have denominators dividing $(q+1)!Q^{q+1}$ and bounded
numerators. The inverse norm and its common denominator are at most
$e^{C(q+1)^2\log(q+2)}Q^{C(q+1)^2}$.
The low-moment realization and the higher-moment Riesz restoration
cost at most the same type of bound times
$(M+1)^{2q+3}Q^{q+3/2}$.
Choose the one rounding denominator larger than their product, the
original common low-moment denominator, and the reciprocal desired
positivity margin. All factors have permitted size. The original
density is retained and only these small corrections are made, so no
coefficient denominator enters this rounding precision. This proves
the claim.
\end{proof}

\begin{theorem}[A uniform count bound]\label{uniform}
There is an absolute constant $C$ such that, for all integers
$n>k\ge1$, a permutation-fair collection of $n$ equiprobable finite
dice exists in which $k$ specified dice all have the same number $S$
of faces satisfying
\[
 S\le (n+2^k+2)^{\exp(Ck^2)}.
\]
The other $n-k$ dice may be taken equal-sized, with no asserted bound
on that common size.
\end{theorem}
\begin{proof}
Use the reviewed induction in
\path{../independent_directions/fixed_number_polynomial_dice.tex},
with reserves
$q_j=2^{k-j+1}-1$ and ordered target degree $M=n+2^k$.
The initial one-uniform model has low moments $1/(a+1)$ through
$q_1=2^k-1$, whose common denominator divides $(2^k)!$. Thus its
complexity is at most
$P_1\le(n+2^k+2)^{C(2^k+1)}$.
The extra factor $2^k+1$ in $\log P_1$ is absorbed by the final
$\exp(Ck^2)$ estimate.

Suppose the current pure uniform-grid model has complexity $P_j$ and
reserve $2q+1$. On a cell $[A/Q,(A+1)/Q]$ of old mass $\gamma$,
choose one global $N$ divisible by the common raw-moment denominator,
and put $K=N\gamma$. A choice
$N\le P_j^{C(q+2)}$, made sufficiently large, supplies the local
separated quadratures in Lemma~\ref{local} for every cell.
For every $a\le2q+1$, the normalized local moment satisfies
\[
 K\mu_a=N\sum_{l=0}^a\binom alQ^l(-A)^{a-l}
          \int_{\rm cell} x^l\,d\nu(x)\in\mathbb Z.
\]
Thus the local low-moment targets have common denominator one before
the harmless factors $1/(s+1)$ in the feature targets.
Apply Lemma~\ref{local} simultaneously to all cells with its shared
control matrix and shared grids. After affine rescaling, common low
moment denominators multiply only the common integers $N,L,E,Q^q$;
there is no least common multiple over cells. All these data have
complexity $P_j^{(q+2)^C}$.
Lemma~\ref{reset} then restores one uniform pure grid, with the same
form of bound after enlarging $C$. In particular,
\[
 \log P_{j+1}\le(q_j+2)^C\log P_j.
\]
As $q_j+2\le2^{k-j+2}$, multiplying these factors over $j<k$ gives
\[
 \log P_k\le \exp(C'k^2)\log(n+2^k+2).
\]

Final separated quadrature on every pure cell, with one common global
count and at least $(n-k)^2$ nodes per cluster, costs a fixed further
power of $P_k$. The uniform quadrature proof has fixed polynomial
losses in the numerical density bounds and degrees. It preserves all
rational cluster moments and produces distinct interior nodes.
The reviewed clustered lifting theorem then rationalizes these $k$
finite special dice and corrects the other $n-k$ continuous uniforms.
It preserves the special counts; its possibly enormous denominators
affect only the final equal-sized old dice. This proves the stated
bound on $S$.
\end{proof}

\begin{theorem}[A growing number of subexponential dice]
There is a constant $c>0$ such that, with
$k=\lfloor c\sqrt{\log n}\rfloor$, one may prescribe $k$ dice in a
fully permutation-fair $n$-die family to have the same count
$S\le\exp(n^{1/2})$ for all sufficiently large $n$.
More generally, if $k(n)=o(\sqrt{\log n})$ and $k(n)\to\infty$,
then the same construction gives $S\le\exp(n^{o(1)})$.
\end{theorem}
\begin{proof}
Choose $c$ small enough that $Cc^2<1/3$ in Theorem~\ref{uniform}.
Then $2^k=n^{o(1)}$, and
$\log S\le n^{1/3+o(1)}\log n\le n^{1/2}$.
For the second statement, $\exp(Ck^2)=n^{o(1)}$ and
$\log(n+2^k+2)=n^{o(1)}$.
\end{proof}

The assertion controls only the specified dice. It is not an
$\exp(O(n))$ bound for the total family and does not establish
polynomial equalization. The absolute constants are deliberately
unoptimized; the uniform dependence, rather than a useful numerical
threshold, is the conclusion.
\end{document}
